\documentclass[aspectratio=169,11pt]{beamer}

% ============================================
% PACKAGES
% ============================================
\usepackage[utf8]{inputenc}
\usepackage[T1]{fontenc}
\usepackage[ngerman]{babel}
\usepackage{lmodern}
\usepackage{tcolorbox}
\tcbuselibrary{skins, hooks}
\usepackage{listings}
\usepackage{xcolor}
\usepackage{fontawesome5}
\usepackage{tikz}
\usetikzlibrary{shapes, arrows.meta}

% Modern Font
\IfFileExists{FiraSans.sty}{\usepackage[sfdefault]{FiraSans}}{}
\renewcommand{\familydefault}{\sfdefault}

% ============================================
% COLORS & DESIGN SYSTEM
% ============================================
\definecolor{BgColor}{HTML}{FAFAFA}
\definecolor{Primary}{HTML}{2B3A42}
\definecolor{Accent}{HTML}{E84A5F}
\definecolor{Secondary}{HTML}{3F5765}
\definecolor{CodeBg}{HTML}{F0F0F0}
\definecolor{Success}{HTML}{28A745}

\setbeamercolor{background canvas}{bg=BgColor}
\setbeamercolor{title}{fg=Primary}
\setbeamercolor{frametitle}{bg=Primary, fg=white}
\setbeamercolor{structure}{fg=Accent}
\setbeamercolor{normal text}{fg=Primary}
\setbeamercolor{itemize item}{fg=Accent}

\setbeamertemplate{navigation symbols}{}
\setbeamertemplate{footline}{%
  \leavevmode%
  \hbox{%
  \begin{beamercolorbox}[wd=.333333\paperwidth,ht=2.25ex,dp=1ex,center]{author in head/foot}%
    \usebeamerfont{author in head/foot}\tiny Falko Himmel
  \end{beamercolorbox}%
  \begin{beamercolorbox}[wd=.333333\paperwidth,ht=2.25ex,dp=1ex,center]{title in head/foot}%
    \usebeamerfont{title in head/foot}\tiny Grundlagen der Praktischen Informatik
  \end{beamercolorbox}%
  \begin{beamercolorbox}[wd=.333333\paperwidth,ht=2.25ex,dp=1ex,right]{date in head/foot}%
    \usebeamerfont{date in head/foot}\insertframenumber{} / \inserttotalframenumber\hspace*{2ex} 
  \end{beamercolorbox}}%
  \vskip0pt%
}

% ============================================
% CUSTOM BOXES
% ============================================
\tcbset{
    modernbox/.style={
        enhanced,
        boxrule=0pt,
        frame hidden,
        colback=white,
        coltitle=white,
        colbacktitle=Secondary,
        fonttitle=\bfseries,
        drop lifted shadow,
        arc=4pt,
        toptitle=2mm,
        bottomtitle=2mm,
        left=2mm,
        right=2mm
    },
    alertbox/.style={
        modernbox,
        colbacktitle=Accent
    }
}

% ============================================
% CODE LISTINGS
% ============================================
\lstdefinelanguage{Haskell}{
  keywords={class, data, deriving, do, else, if, in, import, let, module, of, then, type, where, case, error},
  keywordstyle=\color{Accent}\bfseries,
  identifierstyle=\color{Primary},
  sensitive=true,
  comment=[l]{--},
  commentstyle=\color{gray}\ttfamily\itshape,
  stringstyle=\color{teal}\ttfamily,
  morestring=[b]"
}

\lstset{
    language=Haskell,
    backgroundcolor=\color{CodeBg},
    basicstyle=\ttfamily\scriptsize,
    breaklines=true,
    frame=single,
    rulecolor=\color{CodeBg},
    frameround=tttt,
    numbers=left,
    numberstyle=\tiny\color{gray},
    xleftmargin=15pt,
    tabsize=4,
    showstringspaces=false
}

% ============================================
% TITLE PAGE
% ============================================
\title{\textbf{Tutorium 11}}
\subtitle{Suchalgorithmen}
\institute{Grundlagen der Praktischen Informatik}
\author{}
\date{\today}

\begin{document}

\begin{frame}[plain]
    \titlepage
\end{frame}

\begin{frame}{Inhaltsverzeichnis}
    \tableofcontents
\end{frame}

% ============================================
% THEORIE
% ============================================
\section{Suchen in Listen}

\begin{frame}[fragile]{Lineare Suche}
    Die einfachste Methode: Wir überprüfen jedes Element der Liste nacheinander.
    
    \begin{tcolorbox}[modernbox, title=Lineare Suche in Haskell]
\begin{lstlisting}
contains :: Eq a => [a] -> a -> Bool
contains [] _ = False
contains (y:ys) x
    | x == y    = True  
    | otherwise = contains ys x
\end{lstlisting}
    \end{tcolorbox}
    
    \begin{itemize}
        \item \textbf{Laufzeit:} $\mathcal{O}(n)$
        \item Wir müssen im \textit{Worst Case} jedes Element genau einmal betrachten.
        \item Rekurrenzgleichung: $T(n) = T(n-1) + \mathcal{O}(1)$
    \end{itemize}
\end{frame}


\begin{frame}{Lineare Suche (Visuell)}
    \begin{center}
    \begin{tikzpicture}[node distance=1.2cm, every node/.style={draw, rectangle, minimum size=0.8cm, font=\large}]
        \node (A) {1};
        \node (B) [right of=A] {4};
        \node (C) [right of=B] {6};
        \node (D) [right of=C, fill=Success!30] {7};
        \node (E) [right of=D] {9};
        
        \draw[->, thick, Accent] (A.south) -- +(0,-0.5) node[below, draw=none] {\tiny 1? Nein.};
        \draw[->, thick, Accent] (B.south) -- +(0,-0.5) node[below, draw=none] {\tiny 2? Nein.};
        \draw[->, thick, Success] (D.south) -- +(0,-0.5) node[below, draw=none] {\tiny Gefunden!};
    \end{tikzpicture}
    
    \vspace{0.5cm}
    \large Im \textbf{Worst Case} müssen wir die komplette Liste ablaufen.
    \end{center}
\end{frame}

\begin{frame}[fragile]{Binäre Suche}
    Wenn die Liste \textbf{bereits sortiert} ist, können wir viel effizienter vorgehen: Wir halbieren den Suchraum bei jedem Schritt!
    
    \begin{tcolorbox}[modernbox, title=Binäre Suche in Haskell]
\begin{lstlisting}
containsBinary :: Ord a => [a] -> a -> Bool
containsBinary xs x = go xs
  where
    go [] = False
    go ys = let mid = length ys `div` 2
                midVal = ys !! mid
            in case compare x midVal of
                 LT -> go (take mid ys)
                 GT -> go (drop (mid + 1) ys)
                 EQ -> True
\end{lstlisting}
    \end{tcolorbox}
\end{frame}


\begin{frame}{Binäre Suche (Visuell)}
    \begin{center}
    Suche nach \textbf{7}:\\
    \vspace{0.3cm}
    \begin{tikzpicture}[every node/.style={draw, rectangle, minimum size=0.6cm, font=\small}]
        % Initial
        \node (A) at (0,0) {1}; \node (B) at (0.6,0) {4}; \node (C) at (1.2,0) {6}; \node (D) at (1.8,0) {7}; \node (E) at (2.4,0) {9};
        \node (F) at (3.0,0) {10}; \node (G) at (3.6,0) {213}; \node (H) at (4.2,0) {534}; \node (I) at (4.8,0) {865}; \node (J) at (5.4,0) {999};
        
        \draw[->, thick] (2.7, -0.4) -- (2.7, -0.8);
        
        % Step 1: Halved
        \node[fill=gray!30] (A2) at (0,-1.5) {1}; \node[fill=gray!30] (B2) at (0.6,-1.5) {4}; \node[fill=gray!30] (C2) at (1.2,-1.5) {6}; \node[fill=gray!30] (D2) at (1.8,-1.5) {7}; \node[fill=gray!30] (E2) at (2.4,-1.5) {9};
        
        \node[draw=none, fill=none] at (4.2, -1.5) {\footnotesize (Rechte Hälfte verworfen)};
        
        \draw[->, thick] (1.2, -1.9) -- (1.2, -2.3);
        
        % Step 2
        \node[fill=gray!30] (C3) at (1.2,-3) {6}; \node[fill=Success!30] (D3) at (1.8,-3) {7}; \node[fill=gray!30] (E3) at (2.4,-3) {9};
        
    \end{tikzpicture}
    
    \vspace{0.5cm}
    \large Bei jedem Schritt halbieren wir den Suchraum! $\rightarrow \mathcal{O}(\log n)$
    \end{center}
\end{frame}

\begin{frame}{Effizienz der Binären Suche}
    \begin{itemize}
        \item \textbf{Laufzeit:} $\mathcal{O}(\log_2 n)$
        \item Da die Liste in jedem Schritt halbiert wird, entspricht die Anzahl der Schritte der Höhe eines Binärbaums.
        \item \textbf{Beispiel:} Bei einer Million Einträge benötigen wir \textbf{maximal 20 Vergleiche}!
    \end{itemize}
    
    \vspace{0.5cm}
    \begin{center}
        \texttt{[1, 4, 6, 7, 9, 10, 213, 534, 865, 999]}
    \end{center}
    Wir vergleichen in der Mitte und werfen jeweils die unpassende Hälfte weg.
\end{frame}

% ============================================
% LIVE DEMO
% ============================================
\section{Praxis}

\begin{frame}[plain]
    \begin{center}
        \Huge \textbf{Live Demo} \\
        \vspace{0.5cm}
        \large \faLaptopCode~ \texttt{05\_Suchen.txt} \\
        \vspace{0.3cm}
        \normalsize Gemeinsame Analyse und Programmierung der Suchalgorithmen.
    \end{center}
\end{frame}

\end{document}
